% Fuente única de clase: Recursos/2025-10-30.pdf, páginas 1--13.
En la clase 19 relacionamos raíces e irreducibilidad, construimos extensiones que contienen una raíz y caracterizamos los campos algebraicamente cerrados. Estudiamos después los polinomios reales y la reducción módulo un primo para los polinomios enteros. Terminamos con una forma general del criterio de Eisenstein.

%%%%%%%%%%%%%%%%%%%%%%%%%%%%%%%%%%%%%%%%%%%%%%%%%%%%%%%%%%%%%%%%%%%%%%%%%%%%%%%
\section{Raíces, multiplicidades e irreducibilidad}
%%%%%%%%%%%%%%%%%%%%%%%%%%%%%%%%%%%%%%%%%%%%%%%%%%%%%%%%%%%%%%%%%%%%%%%%%%%%%%%
%%%%%%%%%%%%%%%%%%%%%%%%%%%%%%%%%%%%%%%%%%%%%%%%%%%%%%%%%%%%%%%%%%%%%%%%%%%%%%%
\subsection{El número de raíces y la evaluación}
%%%%%%%%%%%%%%%%%%%%%%%%%%%%%%%%%%%%%%%%%%%%%%%%%%%%%%%%%%%%%%%%%%%%%%%%%%%%%%%
% P01--P02.
En esta clase los anillos son conmutativos. Un elemento $a\in R$ es una \emph{raíz} de $f\in R[x]$ si $f(a)=0$. Por la división por $x-a$, esto equivale a $x-a\mid f$. Para $f\ne0$, la raíz tiene \emph{multiplicidad} $r$ si $(x-a)^r\mid f$ pero $(x-a)^{r+1}\nmid f$.

Las gráficas ilustran perfiles locales de raíces reales de multiplicidades uno, dos y tres.
% Tres perfiles de 2025-10-30.pdf P01: raíz simple, doble y triple en cero.
% Ejemplos ilustrativos; no son gráficas de polinomios especificados allí.

\begin{tikzpicture}[x=1.35cm,y=1.15cm,line width=.6pt,font=\small]
\foreach \x/\y/\m in {0/0/1,5/0/2,2.5/-3.5/3}{
  \begin{scope}[shift={(\x,\y)}]
  \draw[grisCIMAT] (-.9,0)--(2.05,0);
  \draw[grisCIMAT] (0,-1)--(0,1.35);
  \node[below] at (.65,-1.15) {Multiplicidad $\m$};
  \end{scope}
}
\begin{scope}
\clip (-.9,-1) rectangle (2.05,1.35);
\draw[black,samples=100,domain=-.9:2.05,smooth] plot (\x,{.8*\x*(\x-1.6)});
\end{scope}
\begin{scope}[shift={(5,0)}]
\clip (-.9,-1) rectangle (2.05,1.35);
\draw[black,samples=100,domain=-.9:2.05,smooth] plot (\x,{3*\x*\x*(1.2-\x)});
\end{scope}
\begin{scope}[shift={(2.5,-3.5)}]
\clip (-.9,-1) rectangle (2.05,1.35);
\draw[black,samples=100,domain=-.9:2.05,smooth] plot (\x,{\x*\x*\x*(\x-1.5)});
\end{scope}
\end{tikzpicture}


\begin{lema}\label{lema:C19-numero-raices}
Un polinomio no nulo de grado $n$ sobre un dominio tiene a lo sumo $n$ raíces, contadas con multiplicidad.
\end{lema}
\begin{proof}
Sea $K$ el campo de fracciones del dominio. Las raíces del anillo original siguen siendo raíces en $K$, y sus multiplicidades se conservan: al dividir por el mónico $(x-a)^r$, el resto calculado sobre el anillo es también el resto sobre $K$, y la inclusión de coeficientes es inyectiva.

En $K[x]$ hay factorización única y las raíces corresponden a los factores irreducibles de grado uno. Los factores $x-a$ de raíces distintas no son asociados. La suma de sus multiplicidades no puede exceder la suma de los grados de todos los factores de $f$, que es $n$.
\end{proof}
\begin{ejem}
El polinomio $2x-1$ no tiene raíces en $\mathbb Z$, pero sí en $\mathbb Q$, donde su raíz es $1/2$. El paso al campo de fracciones puede añadir raíces.

La condición de dominio es necesaria: en $(\mathbb Z/6\mathbb Z)[x]$, el polinomio $x^2+x$ tiene las cuatro raíces $0,2,3,5$, pues los productos $a(a+1)$ correspondientes son $0,6,12,30$ y todos se anulan módulo seis.
\end{ejem}

\begin{coro}\label{coro:C19-funciones}
Sobre un dominio infinito, dos polinomios son iguales si y sólo si inducen la misma función de evaluación.
\end{coro}
\begin{proof}
Una igualdad formal da iguales evaluaciones. A la inversa, si $f(a)=g(a)$ para todo $a$, cada elemento del dominio es raíz de $f-g$. Si esta diferencia fuera no nula, tendría infinitas raíces, contra el lema.
\end{proof}
Para cada $r\in R$, la evaluación es el único homomorfismo de $R$-álgebras $\varphi_r:R[x]\to R$ que envía $x$ a $r$; en particular,
\[
\varphi_r\left(\sum_i a_ix^i\right)=\sum_i a_ir^i.
\]
Así, el polinomio formal $f$ determina la función $\widehat f(r)=\varphi_r(f)$. La evaluación se representa mediante el triángulo
\[
\begin{tikzcd}[column sep=large,row sep=large]
R[x]\arrow[r,"\varphi_r"]&R\\
\{r\}\arrow[u,"j"]\arrow[ur,hook]&
\end{tikzcd}
\qquad j(r)=x.
\]
La unicidad requiere fijar los coeficientes de $R$ y la imagen de $x$, como en la propiedad universal de la clase 14. El \cref{ejer:C18-formal-funcion} muestra por qué la infinitud importa.

%%%%%%%%%%%%%%%%%%%%%%%%%%%%%%%%%%%%%%%%%%%%%%%%%%%%%%%%%%%%%%%%%%%%%%%%%%%%%%%
\subsection{Grados dos y tres y un campo de cuatro elementos}
%%%%%%%%%%%%%%%%%%%%%%%%%%%%%%%%%%%%%%%%%%%%%%%%%%%%%%%%%%%%%%%%%%%%%%%%%%%%%%%
% P03--P04.
\begin{prop}\label{prop:C19-grados-bajos}
Sobre un campo $k$, un polinomio de grado dos o tres es irreducible si y sólo si no tiene raíces en $k$.
\end{prop}
\begin{proof}
Si $a$ es raíz, escribimos $f=(x-a)g$. Como $\deg g=\deg f-1\geq1$, ninguno de los factores es unidad, así que $f$ es reducible.

Recíprocamente, si $f=gh$ con ambos factores no unidades, sus grados son positivos y suman dos o tres. Alguno tiene grado uno, digamos $g=ux+v$ con $u\ne0$. Como $k$ es campo, $-u^{-1}v$ pertenece a $k$ y es raíz de $g$, y por tanto de $f$.
\end{proof}
Más generalmente, un irreducible de grado al menos dos sobre un campo no puede tener una raíz. La vuelta anterior depende tanto del grado como de la hipótesis de campo. Por ejemplo, $x^2+1$ es irreducible en $\mathbb R[x]$ porque no tiene raíces reales.

\begin{ejem}
En $\mathbb F_2[t]$, el polinomio $t^2+t+1$ toma el valor $1$ en $0$ y en $1$. Es irreducible, su ideal es maximal y
\[
F=\mathbb F_2[t]/(t^2+t+1)
\]
es un campo con cuatro elementos. Si $\alpha$ es la clase de $t$, cada elemento se escribe de manera única $a+b\alpha$, con $a,b\in\mathbb F_2$. Son $0,1,\alpha,1+\alpha$, y $\alpha^2=\alpha+1$. En particular, $\alpha(\alpha+1)=1$.
\end{ejem}
\begin{obs}
La identificación de $F$ con $\mathbb F_2^{\oplus2}$ mediante los coeficientes del resto es un isomorfismo de espacios vectoriales. El producto coordenado de anillos es distinto: $F$ es campo, mientras que $(1,0)(0,1)=(0,0)$ en ese producto. Para multiplicar restos hay que reducir de nuevo usando $\alpha^2=\alpha+1$.
\end{obs}

\begin{ejem}
Si $f\in\mathbb Z[x]$ es primitivo y de grado dos o tres, el \cref{coro:C18-criterio-irred} y la proposición dan
\[
f\text{ irreducible en }\mathbb Z[x]
\quad\Longleftrightarrow\quad
f\text{ no tiene raíces en }\mathbb Q.
\]
No basta buscar raíces enteras: $4x^2-1=(2x-1)(2x+1)$ es reducible en $\mathbb Z[x]$ y no tiene raíces enteras, aunque sí las racionales $\pm1/2$.

En grado cuatro, $x^4+2x^2+1=(x^2+1)^2$ es reducible en $\mathbb Q[x]$ y no tiene raíces racionales. Por tanto, en grado cuatro la ausencia de raíces ya no implica irreducibilidad.
\end{ejem}

%%%%%%%%%%%%%%%%%%%%%%%%%%%%%%%%%%%%%%%%%%%%%%%%%%%%%%%%%%%%%%%%%%%%%%%%%%%%%%%
\subsection{El criterio de raíces racionales}
%%%%%%%%%%%%%%%%%%%%%%%%%%%%%%%%%%%%%%%%%%%%%%%%%%%%%%%%%%%%%%%%%%%%%%%%%%%%%%%
% P04--P05.
\begin{prop}\label{prop:C19-raices-racionales}
Sea $R$ un DFU, $K=K(R)$ y $f=a_0+a_1x+\cdots+a_nx^n\in R[x]$ no nulo. Si $p/q\in K$ es raíz, con $q\ne0$ y $\mcd(p,q)$ unidad, entonces $p\mid a_0$ y $q\mid a_n$.
\end{prop}
\begin{proof}
Si $p=0$, se tiene $a_0=0$ y $\mcd(0,q)$ unidad implica que $q$ es unidad; ambas divisibilidades se cumplen. Supongamos $p\ne0$. Multiplicando $f(p/q)=0$ por $q^n$, obtenemos
\[
a_0q^n+a_1pq^{n-1}+\cdots+a_np^n=0.
\]
Si $n=0$, esto contradiría $f\ne0$, así que $n\geq1$. Reagrupamos como
\[
\begin{aligned}
a_0q^n&=-p(a_1q^{n-1}+\cdots+a_np^{n-1}),\\
a_np^n&=-q(a_0q^{n-1}+a_1pq^{n-2}+\cdots+a_{n-1}p^{n-1}).
\end{aligned}
\]
Así, $p\mid a_0q^n$ y $q\mid a_np^n$. Como $p,q$ no comparten factores irreducibles, comparar las multiplicidades en el DFU da $p\mid a_0$ y $q\mid a_n$. Si uno de esos coeficientes es cero, su divisibilidad es inmediata.
\end{proof}

\begin{ejem}
Para $f=3-2x+3x^2-2x^3+3x^4-2x^5$, los posibles numeradores de una raíz racional reducida son $\pm1,\pm3$, y los denominadores son $\pm1,\pm2$. La identidad
\[
f=(3-2x)(1+x^2+x^4)
\]
muestra que $3/2$ es la única raíz racional del polinomio, pues $1+x^2+x^4>0$ para todo racional $x$.
\end{ejem}

%%%%%%%%%%%%%%%%%%%%%%%%%%%%%%%%%%%%%%%%%%%%%%%%%%%%%%%%%%%%%%%%%%%%%%%%%%%%%%%
\section{Adjuntar raíces y cierre algebraico}
%%%%%%%%%%%%%%%%%%%%%%%%%%%%%%%%%%%%%%%%%%%%%%%%%%%%%%%%%%%%%%%%%%%%%%%%%%%%%%%
%%%%%%%%%%%%%%%%%%%%%%%%%%%%%%%%%%%%%%%%%%%%%%%%%%%%%%%%%%%%%%%%%%%%%%%%%%%%%%%
\subsection{Extensiones mediante un cociente}
%%%%%%%%%%%%%%%%%%%%%%%%%%%%%%%%%%%%%%%%%%%%%%%%%%%%%%%%%%%%%%%%%%%%%%%%%%%%%%%
% P06--P08.
Una \emph{extensión de campos} es un homomorfismo unitario $i:k\to F$. Por el \cref{ejer:C18-morfismo-campos}, es necesariamente inyectivo. Tras identificar $k$ con su imagen escribimos $k\subseteq F$.

\begin{prop}\label{prop:C19-adjuntar}
Sean $k$ un campo y $f(t)\in k[t]$ irreducible. Entonces $F=k[t]/(f)$ es una extensión de $k$, y la clase $\alpha$ de $t$ es una raíz de $f$ en $F$. Para cualquier extensión $j:k\hookrightarrow K$ y raíz $u\in K$ de $f$, existe un único homomorfismo $\widetilde\varphi_u:F\hookrightarrow K$ que extiende $j$ y envía $\alpha$ a $u$.
\end{prop}
\begin{proof}
Como $k[t]$ es DIP, el ideal de un irreducible es maximal; por tanto, $F$ es campo. La composición $k\hookrightarrow k[t]\twoheadrightarrow F$ es unitaria y es inyectiva por ser un homomorfismo entre campos. Aplicando la proyección a los coeficientes y a las potencias de $t$,
\[
f(\alpha)=a_0+a_1\alpha+\cdots+a_n\alpha^n
=\overline{f(t)}=0.
\]
Los coeficientes se entienden mediante la inclusión de $k$.

La propiedad universal de polinomios da la evaluación
\[
\varphi_u:k[t]\longrightarrow K,\qquad
\sum_i a_it^i\longmapsto\sum_i j(a_i)u^i.
\]
Como $f(u)=0$, $(f)\subseteq\ker\varphi_u$. La propiedad universal del cociente induce el único homomorfismo $\widetilde\varphi_u$ que cumple el diagrama:
\[
\begin{tikzcd}[column sep=large,row sep=large]
k[t]\arrow[r,"\varphi_u"]\arrow[d,two heads,"\pi"']&K\\
F\arrow[ur,dashed,destacado,"\exists!\,\widetilde\varphi_u"']&
\end{tikzcd}
\qquad
\begin{tikzcd}[column sep=large,row sep=large]
k\arrow[r,hook,"j"]\arrow[d,hook,"i"']&K\\
F\arrow[ur,hook,"\widetilde\varphi_u"']&
\end{tikzcd}
\]
Es inyectivo por ser un homomorfismo entre campos. La elección de $u$ determina los valores sobre todas las clases de polinomios, de donde resulta la unicidad con esa imagen de $\alpha$.
\end{proof}
\begin{obs}
La evaluación $\varphi_u$ queda determinada al fijar los coeficientes de $k$ y la imagen $t\mapsto u$. El morfismo inducido desde $F$ también depende de la raíz elegida, y raíces distintas pueden dar morfismos distintos.
\end{obs}

\begin{ejem}
Para $k=\mathbb R$ y $f=t^2+1$, se recupera $\mathbb C$ con la construcción de la clase 13. Para $k=\mathbb Q$ y $f=t^2-2$, el polinomio es irreducible: una raíz reducida tendría denominador que divide a $1$ y numerador que divide a $2$, y ninguno de $\pm1,\pm2$ es raíz. El cociente es, pues, un campo con una raíz de $t^2-2$.

Hay exactamente dos inclusiones de ese cociente en $\mathbb R$ que fijan $\mathbb Q$, dadas por
\[
\overline{a+bt}\longmapsto a+b\sqrt2,
\qquad
\overline{a+bt}\longmapsto a-b\sqrt2.
\]
Son distintas y cualquier otra debe enviar la clase de $t$ a una de las dos raíces reales. Esto ilustra por qué no hay unicidad sin elegir la raíz.
\end{ejem}

%%%%%%%%%%%%%%%%%%%%%%%%%%%%%%%%%%%%%%%%%%%%%%%%%%%%%%%%%%%%%%%%%%%%%%%%%%%%%%%
\subsection{Campos algebraicamente cerrados}
%%%%%%%%%%%%%%%%%%%%%%%%%%%%%%%%%%%%%%%%%%%%%%%%%%%%%%%%%%%%%%%%%%%%%%%%%%%%%%%
% P08--P09.
\begin{defi}
Un campo $k$ es \emph{algebraicamente cerrado} si todos los polinomios irreducibles de $k[x]$ tienen grado uno.
\end{defi}
\begin{lema}\label{lema:C19-cerrado}
Para un campo $k$, son equivalentes:
\begin{enumerate}[label=\textup{(\roman*)}]
\item $k$ es algebraicamente cerrado.
\item Todo polinomio no constante de $k[x]$ es producto de factores lineales.
\item Todo polinomio no constante de $k[x]$ tiene una raíz en $k$.
\end{enumerate}
\end{lema}
\begin{proof}
Para (i)$\Rightarrow$(ii), $k[x]$ es DFU, así que un polinomio no constante es producto de irreducibles. Por (i), cada uno tiene grado uno. Reuniendo los coeficientes principales se escribe $f=a\prod_{i=1}^n(x-r_i)$, con $a\in k^\times$.

Para (ii)$\Rightarrow$(iii), hay al menos un factor lineal porque el grado es positivo, y su raíz es raíz de $f$.

Para (iii)$\Rightarrow$(i), sea $f$ irreducible. No es constante, pues las constantes no nulas son unidades y el cero no es irreducible. Por (iii), tiene una raíz $r$, de modo que $f=(x-r)g$. La irreducibilidad hace que $g$ sea unidad, luego constante no nula, y $\deg f=1$.
\end{proof}

\begin{prop}\label{prop:C19-infinito}
Todo campo algebraicamente cerrado es infinito.
\end{prop}
\begin{proof}
Si $k=\{c_1,\ldots,c_N\}$ fuera finito, el polinomio
\[
P(x)=\prod_{i=1}^N(x-c_i)+1
\]
tendría grado $N>0$ y satisfaría $P(c)=1\ne0$ para todo $c\in k$. Contradice el \cref{lema:C19-cerrado}.
\end{proof}

%%%%%%%%%%%%%%%%%%%%%%%%%%%%%%%%%%%%%%%%%%%%%%%%%%%%%%%%%%%%%%%%%%%%%%%%%%%%%%%
\subsection{Polinomios sobre los complejos y los reales}
%%%%%%%%%%%%%%%%%%%%%%%%%%%%%%%%%%%%%%%%%%%%%%%%%%%%%%%%%%%%%%%%%%%%%%%%%%%%%%%
% P10--P11.
\begin{teo}[Teorema fundamental del álgebra]\label{teo:C19-TFA}
El campo $\mathbb C$ es algebraicamente cerrado.
\end{teo}
Usaremos este teorema sin demostración.

\begin{prop}\label{prop:C19-reales}
Los irreducibles de $\mathbb R[x]$ son los polinomios de grado uno y los cuadráticos $ax^2+bx+c$ con $b^2-4ac<0$. Todo polinomio real de grado al menos tres es reducible.
\end{prop}
\begin{proof}
Sea $f\in\mathbb R[x]$ no constante. Por el teorema fundamental del álgebra tiene una raíz compleja $z$. Como los coeficientes son reales,
\[
f(\overline z)=\overline{f(z)}=0.
\]
Si $z\in\mathbb R$, $x-z$ es un factor real. Si $z\ne\overline z$, los factores lineales $x-z$ y $x-\overline z$ son irreducibles no asociados en $\mathbb C[x]$, luego su producto divide a $f$ allí. Ese producto es
\[
q=x^2-(z+\overline z)x+z\overline z\in\mathbb R[x].
\]
También divide a $f$ en $\mathbb R[x]$: al dividir por el mónico $q$ sobre $\mathbb R$, el resto es el mismo que sobre $\mathbb C$ y, por unicidad, es cero. Este paso completa la justificación de que el factor obtenido es real.

Así, todo polinomio no constante tiene un factor real de grado uno o dos. Si su grado es al menos tres, ese factor y el cociente tienen grados positivos, luego es reducible. Los de grado uno son irreducibles. Por la \cref{prop:C19-grados-bajos}, un cuadrático es irreducible exactamente cuando no tiene raíces reales; por la fórmula cuadrática, esto equivale a discriminante negativo.
\end{proof}

\begin{ejer}[Grado impar y una raíz real]\label{ejer:C19-impar}
Probar que todo polinomio real de grado impar tiene una raíz real. Dar una demostración algebraica y otra mediante el teorema del valor intermedio.
\end{ejer}
\begin{sol}
Factoricemos en irreducibles reales. Por la proposición, sus grados son uno o dos. Si todos fueran cuadráticos, el grado total sería par. Luego hay un factor lineal, cuya raíz también anula al polinomio original.

Para el argumento analítico, escribamos $f(x)=a_nx^n+\cdots+a_0$, con $n$ impar y $a_n\ne0$. Para $|x|$ suficientemente grande, $f(x)/(a_nx^n)$ está tan cerca de $1$ como se desee, pues los cocientes de los términos inferiores tienden a cero. Así, $f$ toma signos opuestos en un número positivo suficientemente grande y en su opuesto. La continuidad y el teorema del valor intermedio producen una raíz entre ambos.
\end{sol}

%%%%%%%%%%%%%%%%%%%%%%%%%%%%%%%%%%%%%%%%%%%%%%%%%%%%%%%%%%%%%%%%%%%%%%%%%%%%%%%
\section{Reducción modular y criterio de Eisenstein}
%%%%%%%%%%%%%%%%%%%%%%%%%%%%%%%%%%%%%%%%%%%%%%%%%%%%%%%%%%%%%%%%%%%%%%%%%%%%%%%
%%%%%%%%%%%%%%%%%%%%%%%%%%%%%%%%%%%%%%%%%%%%%%%%%%%%%%%%%%%%%%%%%%%%%%%%%%%%%%%
\subsection{Detectar irreducibilidad módulo un primo}
%%%%%%%%%%%%%%%%%%%%%%%%%%%%%%%%%%%%%%%%%%%%%%%%%%%%%%%%%%%%%%%%%%%%%%%%%%%%%%%
% P11.
La proyección $\mathbb Z\to\mathbb F_p$ induce el homomorfismo sobreyectivo
\[
\mathbb Z[x]\longrightarrow\mathbb F_p[x],\qquad
f\longmapsto\overline f=f\bmod p,
\]
que reduce los coeficientes. Es el mapa del \cref{ejer:C17-coeficientes}.

\begin{prop}\label{prop:C19-modular}
Sea $f\in\mathbb Z[x]$ primitivo y $p$ un primo. Si $\overline f$ tiene el mismo grado que $f$ y es irreducible en $\mathbb F_p[x]$, entonces $f$ es irreducible en $\mathbb Z[x]$ y en $\mathbb Q[x]$.
\end{prop}
\begin{proof}
Si $f=gh$ fuera una factorización en no unidades de $\mathbb Z[x]$, la primitividad y Gauss impedirían que alguno de los factores fuera una constante no unidad. Por tanto, $d=\deg g$ y $e=\deg h$ serían positivos y sumarían $\deg f$. Como $\overline f=\overline g\,\overline h$ conserva el grado,
\[
d+e=\deg\overline f=\deg\overline g+\deg\overline h,
\qquad \deg\overline g\leq d,\quad\deg\overline h\leq e.
\]
Las dos desigualdades deben ser igualdades. Ambos factores reducidos tienen grado positivo, contradiciendo la irreducibilidad de $\overline f$. La afirmación sobre $\mathbb Q$ sigue del \cref{coro:C18-criterio-irred}.
\end{proof}
\begin{ejem}
El polinomio $x^2+x+1$ es primitivo y su reducción módulo dos es irreducible de grado dos, como vimos al construir el campo de cuatro elementos. Por tanto, es irreducible sobre $\mathbb Z$ y sobre $\mathbb Q$.

En cambio,
\[
2x^3+3x^2+3x+1=(2x+1)(x^2+x+1)
\]
es reducible, aunque su reducción módulo dos sea $x^2+x+1$. Aquí el grado baja de tres a dos: falta precisamente la hipótesis de conservación del grado.
\end{ejem}

\begin{samepage}
%%%%%%%%%%%%%%%%%%%%%%%%%%%%%%%%%%%%%%%%%%%%%%%%%%%%%%%%%%%%%%%%%%%%%%%%%%%%%%%
\subsection{Irreducibles de grados arbitrariamente grandes}
%%%%%%%%%%%%%%%%%%%%%%%%%%%%%%%%%%%%%%%%%%%%%%%%%%%%%%%%%%%%%%%%%%%%%%%%%%%%%%%
% P12: ejercicio V.5.11 aplicado y demostrado a mano.
\begin{ejer}[Grados de irreducibles sobre un campo finito]\label{ejer:C19-grados-finitos}
Probar que los grados de los polinomios irreducibles sobre un campo finito no están acotados.
\end{ejer}
\end{samepage}
\begin{sol}
Si todos los irreducibles tuvieran grado a lo sumo $N$, habría sólo finitos irreducibles mónicos $f_1,\ldots,f_m$: hay finitas posibilidades para cada coeficiente y finitos grados posibles. La lista no es vacía porque contiene los polinomios lineales mónicos.

El polinomio $P=f_1\cdots f_m+1$ tiene grado positivo. Como $k[x]$ es DFU, tiene un factor irreducible; normalizándolo, sería alguno de los $f_j$. Pero entonces $f_j$ dividiría tanto a $P$ como a $f_1\cdots f_m$, luego dividiría su diferencia $1$. Esto contradice que tenga grado positivo.
\end{sol}
\begin{coro}
Los grados de los irreducibles de $\mathbb Z[x]$ y de $\mathbb Q[x]$ no están acotados.
\end{coro}
\begin{proof}
Fijemos un primo $p$. Para cada irreducible mónico de $\mathbb F_p[x]$, elegimos representantes enteros de los coeficientes conservando el coeficiente principal $1$. El polinomio obtenido es mónico, luego primitivo, y su reducción conserva el grado. La \cref{prop:C19-modular} da irreducibilidad tanto en $\mathbb Z[x]$ como en $\mathbb Q[x]$. Aplicamos el ejercicio a grados sin cota.
\end{proof}
No se afirma aquí que el argumento produzca todos los grados posibles; prueba exactamente la ausencia de una cota superior.

%%%%%%%%%%%%%%%%%%%%%%%%%%%%%%%%%%%%%%%%%%%%%%%%%%%%%%%%%%%%%%%%%%%%%%%%%%%%%%%
\subsection{La forma general de Eisenstein}
%%%%%%%%%%%%%%%%%%%%%%%%%%%%%%%%%%%%%%%%%%%%%%%%%%%%%%%%%%%%%%%%%%%%%%%%%%%%%%%
% P13 sólo enuncia V.5.17. La demostración siguiente es un complemento añadido.
\begin{prop}[Eisenstein, forma general]\label{prop:C19-Eisenstein}
Sea $R$ un anillo conmutativo y $\mathfrak p$ un ideal primo. Si $f=a_0+a_1x+\cdots+a_nx^n$ satisface
\[
a_n\notin\mathfrak p,\qquad
a_i\in\mathfrak p\ (0\leq i<n),\qquad
a_0\notin\mathfrak p^2,
\]
entonces $f$ no es producto de dos polinomios de grados menores que $n$ en $R[x]$.
\end{prop}
\begin{proof}
Supongamos $f=gh$ con $\deg g,\deg h<n$. En el dominio $D=R/\mathfrak p$, la reducción satisface
\[
\overline g\,\overline h=\overline{a_n}x^n\ne0.
\]
En particular, los dos factores reducidos son no nulos. Para un polinomio no nulo sobre un dominio, tanto su mayor como su menor exponente con coeficiente no nulo se suman al multiplicar: los coeficientes extremos son productos no nulos. Por tanto, si un producto es un monomio, ambos factores son monomios. Escribamos $\overline g=bx^r$, $\overline h=cx^s$, con $b,c\ne0$ y $r+s=n$.

Como $r\leq\deg g<n$ y $s\leq\deg h<n$, los dos exponentes son positivos. Así, los términos constantes de $g$ y $h$ pertenecen a $\mathfrak p$. Su producto es $a_0$, que entonces pertenece a $\mathfrak p^2$, contra la hipótesis. No existe la factorización supuesta. Si $n=0$, tampoco existen factores no nulos de grado menor que cero, por lo que el enunciado es inmediato.
\end{proof}
La conclusión sobre grados es válida incluso si $R$ tiene divisores de cero. En un dominio factorial, para un polinomio primitivo de grado positivo, implica irreducibilidad: una factorización en no unidades no puede tener un factor constante no unidad por Gauss, y los grados positivos de ambos factores serían menores que $n$.

\begin{obs}[Referencias de lectura]
Para raíces, extensiones e irreducibilidad, véase \citet[cap.~V, apartados 5.1--4 y resultados 5.1--17]{aluffi2009}. El teorema fundamental del álgebra se usa sin prueba. Los ejercicios de grado impar y campos finitos corresponden a V.5.17 y V.5.11, respectivamente; la proposición V.5.17 es el criterio de Eisenstein.
\end{obs}
